% ============================================================
% EXTERNAL SOURCE REFERENCES
% Viability Philosophy System
% Book I: Ontology of Viability
% ============================================================

\cleardoublepage

\section*{External Source References}

\phantomsection
\addcontentsline{toc}{section}{External Source References}

% ------------------------------------------------------------
% 1. PURPOSE AND SCOPE
% ------------------------------------------------------------

\subsection*{1. Purpose and Scope}

External source references identify and locate
materials originating outside the internal
logical and structural organization of the
Viability Philosophy System.

They provide the means by which external
books, articles, philosophical works,
documents, editions, and other relevant
materials are acknowledged and referenced
within the text.

The purpose of external source referencing
is not merely to indicate that an external
work has been consulted. It is to make
the relationship between the internal
discussion and the external material
transparent, traceable, and verifiable.

External source references must allow
the reader to distinguish between:

\begin{itemize}
    \item an internal logical entity belonging
    to the Viability Philosophy System;
    \item an internal textual location within
    one of its books;
    \item an external work or document cited
    in the discussion;
    \item a specific passage, page, section,
    or other location within that external work.
\end{itemize}

These distinctions are necessary because
the identity of an internal entity, the
location of its discussion, and the
bibliographic identity of an external
source are different forms of identification.

External source references serve the following
purposes:

\begin{itemize}
    \item identifying external works used
    in the development of the discussion;
    \item locating the specific passages
    relevant to a claim or interpretation;
    \item distinguishing direct quotations
    from paraphrases and general references;
    \item making the use of external sources
    transparent;
    \item supporting verification of
    textual and interpretive claims;
    \item maintaining consistency between
    in-text references and bibliographic
    records.
\end{itemize}

% ------------------------------------------------------------
% 2. EXTERNAL SOURCES AND INTERNAL
% IDENTIFICATION
% ------------------------------------------------------------

\subsection*{2. External Sources and Internal
Identification}

External sources do not belong to the
internal structural hierarchy of the
Viability Philosophy System merely
because they are cited or discussed
within one of its books.

They must therefore not receive internal
structural identifiers.

Likewise, external authors, books,
articles, and documents must not
receive internal logical entity
identifiers merely by being included
in the discussion.

For example:

\begin{center}
\begin{tabular}{ll}
    \toprule
    \textbf{Identification} & \textbf{Function} \\
    \midrule
    \texttt{VS.B01.P01.C01.S01.SS01.P001}
        & Internal textual location \\
    \texttt{AX-001}
        & Internal logical entity \\
    Bibliographic reference
        & External work or document \\
    \bottomrule
\end{tabular}
\end{center}

An external work retains its own
bibliographic identity.

Its appearance in a particular
chapter or paragraph does not
transform it into an internal
component of the book.

Similarly, an internal logical
entity may be discussed in relation
to an external source without
becoming identical to that source.

The distinction between internal
identification and external source
identification must be maintained
throughout the system.

% ------------------------------------------------------------
% 3. IDENTIFICATION OF EXTERNAL WORKS
% ------------------------------------------------------------

\subsection*{3. Identification of External Works}

An external source must be identified
through bibliographic information
sufficient to distinguish it from
other works and to enable its
retrieval.

The information required depends
on the type of source and the
bibliographic conventions adopted
for the book.

Relevant information may include:

\begin{itemize}
    \item the author or responsible
    creator;
    \item the title of the work;
    \item the editor, translator,
    or other relevant contributor;
    \item the edition or version;
    \item the publication date;
    \item the publisher or issuing
    institution;
    \item the place of publication,
    where required by the adopted
    bibliographic convention;
    \item the journal, collection,
    or larger work in which the
    material appears;
    \item volume, issue, or other
    publication details;
    \item a persistent identifier
    or stable retrieval location,
    where applicable.
\end{itemize}

Not every source requires every
item of information.

The bibliographic description
must include the information
appropriate to the source type
and sufficient to distinguish
the cited work from other
potentially similar works.

The precise arrangement and
formatting of these elements
must follow the bibliographic
conventions established for
the book.

% ------------------------------------------------------------
% 4. WORK IDENTITY AND EDITION
% ------------------------------------------------------------

\subsection*{4. Work Identity and Edition}

The identity of a work must
be distinguished from the
identity of a particular
edition, translation, or version.

The same work may exist in
multiple editions or translations.
These versions may differ in
pagination, wording, organization,
or other bibliographically
relevant features.

A reference must identify
the version actually consulted
when the distinction is relevant
to the verification of the
claim or passage.

Where a quotation or specific
textual interpretation depends
on a particular translation
or edition, the reference
should make that version
identifiable.

A reference must not silently
combine the bibliographic
identity of one edition with
the textual location of another.

Where several editions are
used, the bibliographic
records and citations must
preserve the distinctions
necessary for accurate
verification.

% ------------------------------------------------------------
% 5. REFERENCES TO SPECIFIC
% PASSAGES
% ------------------------------------------------------------

\subsection*{5. References to Specific Passages}

A reference to an external work
may identify the work as a whole
or direct the reader to a
specific passage within it.

Where a particular claim, quotation,
interpretation, or textual detail
depends on a specific passage,
the reference should provide
an appropriate locator.

Depending on the source, such
a locator may consist of:

\begin{itemize}
    \item a page or page range;
    \item a chapter;
    \item a section or subsection;
    \item a numbered paragraph;
    \item a book, part, or other
    internal division of the work;
    \item a stable section or
    passage identifier;
    \item a timestamp or equivalent
    locator for audiovisual material;
    \item another source-specific
    means of locating the material.
\end{itemize}

The locator must correspond
to the version of the source
actually consulted.

A page reference must not be
treated as universally applicable
when different editions or
translations have different
pagination.

Where pagination is unavailable
or unsuitable, an alternative
locator should be used when
the source provides one.

The purpose of a locator is
to make the cited material
verifiable, not merely to
add numerical detail to a
bibliographic reference.

% ------------------------------------------------------------
% 6. DIRECT QUOTATIONS
% ------------------------------------------------------------

\subsection*{6. Direct Quotations}

A direct quotation reproduces
wording from an external source.

The quoted wording must be
accurate and must preserve
the meaning of the original
passage.

A direct quotation should
be accompanied by a reference
that identifies the source
and, where available and
appropriate, the specific
location of the quoted text.

The following principles apply:

\begin{itemize}
    \item Quoted wording must
    not be presented as the
    author's own formulation.

    \item Omissions must not
    distort the meaning of
    the original passage.

    \item Editorial additions,
    clarifications, or alterations
    within a quotation must
    be clearly distinguished
    from the source's wording.

    \item The source and the
    relevant textual location
    must be identifiable.

    \item The quotation must
    be represented in a way
    that preserves the distinction
    between the original text
    and the surrounding discussion.
\end{itemize}

When a quotation is translated
by the author of the present
work rather than reproduced
from a published translation,
that fact should be made clear
where it is relevant to the
reader's interpretation.

A translated quotation must
not be presented as though
its wording were necessarily
identical to that of an
established published translation.

% ------------------------------------------------------------
% 7. PARAPHRASES AND INTERPRETIVE
% REFERENCES
% ------------------------------------------------------------

\subsection*{7. Paraphrases and Interpretive
References}

A paraphrase presents the
substance of an external
source in different wording.

Although it does not reproduce
the original wording, it remains
dependent on the source and
must be appropriately referenced.

A paraphrase must preserve
the relevant meaning of the
source without presenting
an interpretation as though
it were an exact quotation.

The author must distinguish
between:

\begin{itemize}
    \item what the external
    source explicitly states;
    \item what may reasonably
    be inferred from the source;
    \item the author's own
    interpretation of the source;
    \item the author's critical
    assessment or disagreement.
\end{itemize}

These distinctions are particularly
important when an external
philosophical position is
reconstructed, interpreted,
or critically evaluated.

A reference to a source does
not establish that every
interpretation associated
with that source is explicitly
endorsed by its author.

The wording of the discussion
must therefore make clear
whether it reports, interprets,
reconstructs, or criticizes
the cited position.

% ------------------------------------------------------------
% 8. GENERAL REFERENCES
% ------------------------------------------------------------

\subsection*{8. General References}

A general reference identifies
an external work as relevant
to a discussion without necessarily
directing the reader to one
specific passage.

Such a reference may be appropriate
when acknowledging a broader
theoretical context, identifying
a philosophical tradition, or
indicating a work that has
informed the discussion as
a whole.

However, a general reference
must not be used as a substitute
for a specific locator when
a particular passage or claim
requires verification.

The level of precision must
correspond to the nature of
the claim being made.

% ------------------------------------------------------------
% 9. SECONDARY AND INDIRECT
% REFERENCES
% ------------------------------------------------------------

\subsection*{9. Secondary and Indirect References}

A secondary reference occurs
when an external work is
discussed or cited through
another source rather than
consulted directly.

Where the original source
has not been consulted,
the reference must not create
the impression that it has
been examined firsthand.

The distinction between
direct consultation and
indirect reference must
be preserved.

Where possible, the original
source should be consulted
when the argument depends
on its exact wording or
a specific interpretation.

If only a secondary source
is available, the citation
must make the indirect
nature of the reference
clear according to the
adopted bibliographic
conventions.

A secondary source must
not be silently substituted
for the original source
as though the two were
interchangeable.

% ------------------------------------------------------------
% 10. EXTERNAL SOURCES AND
% INTERNAL ARGUMENTS
% ------------------------------------------------------------

\subsection*{10. External Sources and Internal
Arguments}

External sources may be used
for different purposes within
the philosophical argument.

They may provide:

\begin{itemize}
    \item historical or conceptual
    context;
    \item a philosophical position
    to be examined;
    \item an argument or distinction
    relevant to the discussion;
    \item evidence for a historical
    or interpretive claim;
    \item a point of comparison;
    \item a position to be criticized
    or rejected;
    \item a formulation that
    requires clarification or
    reconstruction.
\end{itemize}

The function of an external
source within an argument
must be made sufficiently
clear.

The mere citation of an
external author or work
does not establish that
the present book adopts
the source's conclusions.

Likewise, the use of an
external source does not
automatically make the
source's conceptual framework
part of the internal logical
organization of the Viability
Philosophy System.

The author must distinguish
between the source's position
and the position developed
within the present work.

% ------------------------------------------------------------
% 11. SOURCE ACCURACY AND
% VERIFICATION
% ------------------------------------------------------------

\subsection*{11. Source Accuracy and Verification}

External references must
be accurate enough to allow
the reader to identify
and verify the cited material.

The following requirements
apply:

\begin{itemize}
    \item Author names and
    work titles must be
    represented accurately.

    \item Publication information
    must correspond to the
    version actually consulted.

    \item Quoted wording must
    be checked against the
    source.

    \item Specific locators
    must correspond to the
    cited edition or version.

    \item Translations must
    be identified accurately
    where relevant.

    \item Bibliographic details
    must not be invented
    to complete an incomplete
    reference.

    \item Unverified information
    must not be presented
    as though it had been
    confirmed.
\end{itemize}

Where bibliographic information
is unavailable, the missing
information should be handled
according to the adopted
bibliographic conventions
rather than supplied through
speculation.

The accuracy of an external
reference is a substantive
requirement of the argument,
not merely a matter of
typographical presentation.

% ------------------------------------------------------------
% 12. EXTERNAL SOURCES AND
% INTERNAL CROSS-REFERENCES
% ------------------------------------------------------------

\subsection*{12. External Sources and
Internal Cross-References}

External source references
and internal cross-references
perform different functions.

An external source reference
identifies material outside
the Viability Philosophy
System.

An internal cross-reference
directs the reader to another
location or logical entity
within the system.

The two may appear together
when an internal discussion
refers to an external work
and also directs the reader
to another passage within
the book.

In such cases, each reference
must retain its own function.

An external bibliographic
reference must not be
substituted for an internal
structural reference.

Likewise, an internal
structural identifier must
not be treated as a
bibliographic identifier
for an external work.

% ------------------------------------------------------------
% 13. EXTERNAL SOURCES AND
% LOGICAL ENTITY REFERENCES
% ------------------------------------------------------------

\subsection*{13. External Sources and
Logical Entity References}

An internal logical entity
may be examined in relation
to one or more external
sources.

The logical entity identifier
continues to designate
the internal entity.

The external source reference
continues to identify the
external work.

The relation between them
must be expressed through
the argument and the
relevant references.

An external author or
philosophical position
must not be assigned an
internal logical entity
identifier merely because
it is discussed in connection
with an internal entity.

Similarly, an internal
logical entity must not
be identified solely through
the bibliographic identity
of an external source.

% ------------------------------------------------------------
% 14. DIGITAL AND NON-PRINT
% SOURCES
% ------------------------------------------------------------

\subsection*{14. Digital and Non-Print Sources}

External sources are not
limited to printed books
and articles.

The referencing system
may also accommodate
digital documents, online
publications, institutional
materials, audiovisual
works, archival materials,
and other non-print sources.

The bibliographic description
must be appropriate to
the type of material.

Where relevant, it should
identify:

\begin{itemize}
    \item the responsible
    author or institution;
    \item the title or identifying
    description;
    \item the publication
    or release date;
    \item the version or
    update information;
    \item the publisher or
    issuing body;
    \item a stable URL or
    persistent identifier;
    \item the access date,
    where required by the
    adopted convention;
    \item a suitable locator
    within the material.
\end{itemize}

A URL alone should not
be treated as a complete
bibliographic description
when additional identifying
information is available
or required.

For sources that may change
over time, the cited version
or retrieval context should
be identifiable whenever
this is necessary for
verification.

% ------------------------------------------------------------
% 15. CONSISTENCY OF CITATION
% ------------------------------------------------------------

\subsection*{15. Consistency of Citation}

All external source references
must follow a consistent
bibliographic and citation
system throughout the book.

The same work must not
be represented through
incompatible citation
formats without a clear
and justified reason.

Author names, titles,
publication details, edition
information, and locators
must be represented
consistently.

The form of in-text citations
and the organization of
bibliographic entries must
be coordinated.

The precise citation style,
punctuation, ordering,
abbreviations, and formatting
rules are governed by
the bibliographic conventions
established for the book.

This document establishes
the principles of external
source identification and
referencing; it does not
independently prescribe
a complete bibliographic
style.

% ------------------------------------------------------------
% 16. REVISION AND MAINTENANCE
% ------------------------------------------------------------

\subsection*{16. Revision and Maintenance}

External source references
must be reviewed whenever
the manuscript or its
bibliographic records
are revised.

The following cases require
particular attention.

\subsubsection*{16.1. Change of Edition}

When a different edition
or translation is adopted,
the corresponding references
and locators must be reviewed.

A page reference from one
edition must not be retained
without verification when
the cited edition has changed.

\subsubsection*{16.2. Change of Bibliographic Information}

When publication information
is corrected or completed,
the corresponding in-text
references and bibliographic
records must remain consistent.

\subsubsection*{16.3. Removal of a Source}

When an external source
is removed from the discussion,
its bibliographic record
and all associated references
must be reviewed.

A bibliographic entry must
not remain without a relevant
function merely because
it appeared in an earlier
version of the manuscript.

\subsubsection*{16.4. Addition of a Source}

When a new source is introduced,
its bibliographic information
must be recorded accurately
and its use in the text
must be supported by
appropriate references.

\subsubsection*{16.5. Verification of Existing References}

References must be checked
when quotations, paraphrases,
interpretations, or source
locators are revised.

A change in wording may
require a renewed examination
of whether the cited passage
still supports the claim.

% ------------------------------------------------------------
% 17. RELATION TO BIBLIOGRAPHIC
% ORGANIZATION
% ------------------------------------------------------------

\subsection*{17. Relation to Bibliographic
Organization}

External source references
identify and locate the
external materials used
within the text.

Bibliographic organization
governs the systematic
presentation and arrangement
of the corresponding source
records.

These functions are related
but distinct.

The in-text reference directs
the reader to the relevant
external source and, where
necessary, to a specific
passage within it.

The bibliographic record
provides the identifying
publication information
required to distinguish
and retrieve the source.

Both must remain consistent.

The detailed organization
of bibliographic entries
is established in the
document \textit{Bibliographic
Organization}.

% ------------------------------------------------------------
% 18. SUMMARY
% ------------------------------------------------------------

\subsection*{18. Summary}

External source references
identify and locate materials
originating outside the
internal organization of
the Viability Philosophy
System.

They must remain distinct
from structural identifiers
and logical entity identifiers.

A reliable external reference
must identify the source
accurately and provide
an appropriate locator
when a specific passage
or textual claim is involved.

Direct quotations, paraphrases,
interpretive references,
general references, and
indirect references must
be distinguished according
to their actual relation
to the source.

The use of external sources
must be transparent, verifiable,
and consistent with the
bibliographic conventions
adopted for the book.

External source references
therefore establish the
connection between the
internal philosophical
discussion and the external
materials on which it
draws, without incorporating
those materials into the
internal identification
system of the Viability
Philosophy System.

\cleardoublepage

% ============================================================
% END OF EXTERNAL SOURCE REFERENCES
% ============================================================